% =============================================================
% Post-Quantum Secure API Communication: Architectures and
% Migration Strategies for Quantum-Resilient Systems
%
% Author: Subhajit Roy
% ORCID:  0009-0000-4896-0922
% =============================================================

\documentclass[conference]{IEEEtran}

\usepackage{amsmath,amssymb}
\usepackage{graphicx}
\usepackage{cite}
\usepackage{hyperref}
\usepackage{xcolor}
\usepackage{booktabs}
\usepackage{array}
\usepackage{listings}
\usepackage{algorithm}
\usepackage{algpseudocode}
\usepackage{orcidlink}
\usepackage{url}

\hypersetup{
  colorlinks=true,
  linkcolor=blue,
  citecolor=blue,
  urlcolor=blue,
  pdftitle={Post-Quantum Secure API Communication: Architectures and Migration Strategies for Quantum-Resilient Systems},
  pdfauthor={Subhajit Roy}
}

\begin{document}

\title{Post-Quantum Secure API Communication: Architectures and Migration Strategies for Quantum-Resilient Systems}

\author{
  \IEEEauthorblockN{Subhajit Roy\orcidlink{0009-0000-4896-0922}}
  \IEEEauthorblockA{
    Independent Researcher\\
    Email: subhajitroy857@gmail.com\\
    ORCID: \href{https://orcid.org/0009-0000-4896-0922}{0009-0000-4896-0922}
  }
}

\maketitle

\begin{abstract}
The advent of cryptographically relevant quantum computers threatens the
asymmetric primitives -- RSA, Diffie-Hellman, and elliptic-curve cryptography
-- that underpin the confidentiality and integrity of contemporary
Application Programming Interface (API) communication. While large-scale
quantum adversaries do not yet exist, the ``harvest now, decrypt later''
threat model means that data encrypted today under classical schemes may
already be exposed to future cryptanalysis. This paper examines the
practical migration of API security architectures toward post-quantum
cryptography (PQC), focusing on the National Institute of Standards and
Technology (NIST) standardized algorithms ML-KEM (FIPS 203), ML-DSA (FIPS
204), and SLH-DSA (FIPS 205). We propose a layered reference architecture for
quantum-resilient API systems, evaluate hybrid key-exchange strategies for
Transport Layer Security (TLS) 1.3, and analyze the downstream effects of
PQC adoption on certificate chains, JSON Web Tokens (JWT), OAuth 2.0 token
issuance, and API gateway design. We further present a phased migration
framework -- spanning algorithm inventory, hybrid deployment, and
cryptographic agility -- intended to guide practitioners through the
transition without requiring a disruptive, single-step cutover. The paper
concludes with an assessment of open engineering challenges, including
increased key and signature sizes, handshake latency, and the operational
complexity of maintaining dual classical/post-quantum trust chains during
the transition period.
\end{abstract}

\begin{IEEEkeywords}
post-quantum cryptography, API security, hybrid TLS, ML-KEM, ML-DSA,
cryptographic agility, NIST PQC, TLS 1.3, quantum-safe communication,
applied cryptography
\end{IEEEkeywords}

% =============================================================
\section{Introduction}
% =============================================================

Modern distributed systems rely overwhelmingly on API-mediated
communication, secured at the transport layer by TLS and at the application
layer by token-based authentication and authorization schemes such as JWT
and OAuth 2.0. The cryptographic foundations of these mechanisms --
RSA, Diffie-Hellman key exchange, and elliptic-curve variants thereof --
derive their security from the computational hardness of integer
factorization and the discrete logarithm problem. Shor's algorithm
demonstrates that a sufficiently large fault-tolerant quantum computer
solves both problems in polynomial time \cite{shor1997}, which would render
these primitives insecure.

Although current estimates of when a cryptographically relevant quantum
computer (CRQC) will exist vary considerably across the research community
\cite{mosca_estimate}, the asymmetry between attacker and defender timelines
creates an immediate practical concern. An adversary capable of recording
encrypted API traffic today need not possess a quantum computer at the time
of capture; decryption can be deferred until quantum capability matures.
This ``harvest now, decrypt later'' (HNDL) model means that any
long-lived sensitive data transmitted over classically-secured channels
today is already at risk, independent of when quantum hardware
materializes.

In response, NIST finalized its first set of post-quantum cryptographic
standards in 2024: ML-KEM for key encapsulation (FIPS 203, based on
CRYSTALS-Kyber) \cite{nist_fips203,kyber}, ML-DSA for digital signatures
(FIPS 204, based on CRYSTALS-Dilithium) \cite{nist_fips204,dilithium}, and
SLH-DSA for hash-based digital signatures (FIPS 205, based on SPHINCS+)
\cite{nist_fips205,sphincs}. These standards provide the cryptographic
primitives required for quantum-resistant systems, but their integration
into existing API ecosystems -- TLS termination, certificate authorities,
token issuance services, and API gateways -- requires deliberate
architectural planning.

This paper addresses the system-level question that follows standardization:
\emph{how should existing API infrastructure be re-architected and migrated
to incorporate post-quantum primitives without disrupting availability,
interoperability, or backward compatibility?} We do not propose new
cryptographic constructions; rather, we synthesize current standards
guidance, hybrid key-exchange practice, and migration literature into a
practical reference architecture and phased migration framework intended
for engineering teams operating production API systems.

\subsection{Research Questions}

This work is organized around four research questions:

\begin{enumerate}
  \item \textbf{RQ1:} What architectural layers of an API system are
        affected by the transition to post-quantum cryptography, and what
        are the trust boundaries between them?
  \item \textbf{RQ2:} How should hybrid key exchange be deployed in TLS 1.3
        to provide quantum resistance while preserving compatibility with
        classical clients during the transition period?
  \item \textbf{RQ3:} What changes are required in certificate chains, JWT
        issuance, and OAuth 2.0 flows to accommodate post-quantum signature
        schemes, and what are the practical costs of these changes?
  \item \textbf{RQ4:} What migration sequencing minimizes operational risk
        while ensuring cryptographic agility for future algorithm
        transitions?
\end{enumerate}

\subsection{Key Contributions}

\begin{itemize}
  \item A five-layer reference architecture for quantum-resilient API
        systems, delineating trust boundaries between transport,
        certificate, token, gateway, and application layers.
  \item A threat model for API communication under both classical and HNDL
        quantum adversaries, identifying protected assets and attack
        surfaces specific to API ecosystems.
  \item An analysis of hybrid TLS 1.3 key-exchange deployment using ML-KEM
        alongside classical elliptic-curve Diffie-Hellman (ECDHE), following
        IETF hybrid key-exchange guidance \cite{ietf_hybrid_kex}.
  \item A phased migration framework covering classical-to-hybrid TLS
        transition, certificate chain migration, JWT/OAuth 2.0 signature
        migration, and API gateway re-architecture.
  \item A discussion of cryptographic agility as a design principle,
        enabling future algorithm transitions without structural
        re-architecture.
\end{itemize}

% =============================================================
\section{Background}
% =============================================================

\subsection{The Quantum Threat to Asymmetric Cryptography}

RSA and Diffie-Hellman-based key exchange, including its elliptic-curve
variants (ECDH, ECDHE), rely on problems -- integer factorization and the
discrete logarithm problem -- that admit no known efficient classical
solution. Shor's algorithm \cite{shor1997} solves both in polynomial time
on a quantum computer, which would allow an adversary to recover private
keys from public keys and break both confidentiality (key exchange) and
authenticity (digital signatures) for any system relying on these
primitives. Symmetric primitives such as AES are comparatively less
affected; Grover's algorithm provides only a quadratic speedup, which is
typically mitigated by doubling key length.

\subsection{NIST Post-Quantum Standardization}

NIST's PQC standardization project \cite{nist_pqc_project} selected
lattice-based and hash-based candidates as the first standardized
algorithms:

\begin{itemize}
  \item \textbf{ML-KEM (FIPS 203)} \cite{nist_fips203}, derived from
        CRYSTALS-Kyber \cite{kyber}, provides key encapsulation based on the
        Module Learning With Errors (MLWE) problem.
  \item \textbf{ML-DSA (FIPS 204)} \cite{nist_fips204}, derived from
        CRYSTALS-Dilithium \cite{dilithium}, provides digital signatures
        based on a combination of MLWE and the Module Short Integer
        Solution (MSIS) problem.
  \item \textbf{SLH-DSA (FIPS 205)} \cite{nist_fips205}, derived from
        SPHINCS+ \cite{sphincs}, provides a stateless hash-based signature
        scheme with security resting solely on the collision resistance of
        the underlying hash function, serving as a conservative fallback
        should lattice assumptions weaken.
  \item \textbf{FALCON} \cite{falcon}, a NTRU-lattice-based signature
        scheme offering compact signatures, remains under continued
        standardization consideration as a complementary signature
        algorithm.
\end{itemize}

Table~\ref{tab:algo-comparison} summarizes representative key and signature
sizes relative to classical equivalents at comparable security levels.

\begin{table}[h]
\centering
\caption{Approximate Key and Signature Sizes (Security Category 1, $\approx$128-bit)}
\label{tab:algo-comparison}
\begin{tabular}{@{}lccc@{}}
\toprule
\textbf{Scheme} & \textbf{Public Key} & \textbf{Private Key} & \textbf{Ciphertext/Sig} \\
\midrule
ECDH (X25519)   & 32 B    & 32 B    & 32 B \\
ML-KEM-512      & 800 B   & 1632 B  & 768 B \\
ECDSA (P-256)   & 64 B    & 32 B    & 64--72 B \\
ML-DSA-44       & 1312 B  & 2560 B  & 2420 B \\
SLH-DSA-128s    & 32 B    & 64 B    & 7856 B \\
\bottomrule
\end{tabular}
\end{table}

The order-of-magnitude increase in key and signature sizes, particularly for
SLH-DSA, has direct implications for TLS handshake size, certificate chain
length, and JWT payload size -- each discussed in Section~\ref{sec:arch}.

\subsection{Hybrid Cryptography as a Transition Strategy}

Rather than replacing classical primitives outright, current IETF and NSA
guidance favors \emph{hybrid} key exchange during the transition period
\cite{ietf_hybrid_kex,nsa_cnsa2}. A hybrid key exchange combines a classical
algorithm (e.g., X25519) with a post-quantum algorithm (e.g., ML-KEM) such
that the resulting shared secret is secure if \emph{either} component
remains unbroken. This approach hedges against two risks simultaneously:
undiscovered weaknesses in newly standardized post-quantum schemes, and the
eventual arrival of a CRQC capable of breaking classical schemes. We adopt
hybrid key exchange as the primary transport-layer mechanism in the
architecture proposed in Section~\ref{sec:arch}.

% =============================================================
\section{Threat Model}
% =============================================================

A complete threat model for this work is maintained in
\texttt{docs/threat-model.md} within the accompanying repository; this
section summarizes its core elements.

\subsection{Adversary Capabilities}

We consider two adversary classes:

\begin{itemize}
  \item \textbf{Classical network adversary:} capable of passive
        eavesdropping, active man-in-the-middle positioning, and
        certificate or token forgery attempts bounded by classical
        computational limits.
  \item \textbf{Harvest-now, decrypt-later (HNDL) adversary:} passively
        records encrypted traffic and certificate material today, with the
        intent of applying quantum cryptanalysis once a CRQC becomes
        available. This adversary is bounded only by data retention
        capacity, not by present-day computational limits.
\end{itemize}

\subsection{Protected Assets}

The assets in scope are: (i) the confidentiality of data in transit between
API clients and servers, (ii) the integrity and authenticity of API
requests and responses, (iii) the long-term validity of certificate chains
and trust anchors, and (iv) the integrity of bearer tokens (JWT) used for
authentication and authorization.

\subsection{Security Goals}

The architecture in Section~\ref{sec:arch} is designed to achieve: forward
secrecy against both classical and HNDL adversaries at the transport layer;
unforgeability of certificates and tokens against quantum-capable
signature forgery; and backward compatibility with classical-only clients
during the migration period, without weakening the security posture for
quantum-resistant clients.

% =============================================================
\section{Reference Architecture}
\label{sec:arch}
% =============================================================

We propose a five-layer architecture for reasoning about post-quantum
migration in API systems, illustrated in the API architecture diagram
accompanying this repository (\texttt{diagrams/api-architecture.svg}).

\begin{enumerate}
  \item \textbf{Transport Layer:} TLS 1.3 connection establishment,
        including the key-exchange mechanism (classical, post-quantum, or
        hybrid).
  \item \textbf{Certificate Layer:} X.509 certificate issuance and
        validation, including the signature algorithm used by certificate
        authorities and the structure of certificate chains during
        dual-algorithm operation.
  \item \textbf{Token Layer:} Application-level authentication and
        authorization artifacts -- JWTs and OAuth 2.0 access/refresh tokens
        -- and the signature or MAC scheme used to protect them.
  \item \textbf{Gateway Layer:} API gateway and reverse proxy components
        responsible for TLS termination, request routing, and token
        validation, which must support algorithm negotiation and fallback.
  \item \textbf{Application Layer:} Business logic consuming validated
        requests, generally agnostic to the underlying cryptographic
        scheme provided the lower layers expose a stable interface.
\end{enumerate}

A diagrammatic representation of this architecture, the corresponding
hybrid TLS handshake, and the cryptographic dependency graph are provided
in \texttt{diagrams/} and described in \texttt{docs/architecture.md}.

\subsection{Trust Boundaries}

Each layer boundary represents a point where cryptographic material crosses
from one trust domain to another: the transport layer boundary is the TLS
session itself; the certificate layer boundary is the handoff between
certificate authority and relying party; the token layer boundary is the
handoff between identity provider and resource server; and the gateway
boundary mediates between external, potentially hybrid-capable clients and
internal services that may still operate classical-only cryptography during
migration.

\subsection{Cryptographic Agility}

We treat cryptographic agility -- the capacity to substitute algorithms
without redesigning system architecture -- as a first-class design
requirement rather than an afterthought. Practically, this is realized
through: algorithm identifiers negotiated at the protocol level (TLS
cipher suites, JWT \texttt{alg} headers), configuration-driven algorithm
selection in gateway and token-issuance components, and avoidance of
hard-coded key or signature size assumptions in downstream application
code.

% =============================================================
\section{Hybrid TLS for API Transport Security}
% =============================================================

Hybrid key exchange in TLS 1.3 combines an ML-KEM encapsulation with a
classical ECDHE exchange within the \texttt{key\_share} extension,
following the IETF hybrid design draft \cite{ietf_hybrid_kex}. The combined
shared secret is derived via a key-derivation function (KDF) that mixes
both component secrets, ensuring that the session remains secure if either
the classical or the post-quantum component is broken, but not necessarily
if both are independently broken in different ways -- a property generally
referred to as IND-CCA2 security of the combiner.

\subsection{Handshake Size and Latency}

The principal engineering cost of hybrid key exchange is increased
handshake size: ML-KEM-768 key shares are on the order of 1.1--1.2 KB
versus 32 bytes for X25519, increasing the TLS \texttt{ClientHello} and
\texttt{ServerHello} payloads correspondingly. For most API workloads using
persistent or pooled connections, this one-time handshake cost is
amortized over the connection lifetime and has negligible effect on
steady-state throughput; latency-sensitive, connection-per-request
architectures should weight this cost more heavily in capacity planning.

\subsection{Server and Client Negotiation}

Servers supporting hybrid TLS advertise hybrid key-share groups alongside
classical groups, allowing clients without post-quantum support to
negotiate a purely classical handshake. This preserves interoperability
during the transition but means the security guarantee is only as strong
as the weaker party's capability; gateway-layer policy (Section~\ref{sec:arch})
should be used to enforce minimum acceptable key-exchange groups for
sensitive endpoints once client-side support is sufficiently mature.

% =============================================================
\section{Certificate, JWT, and OAuth 2.0 Migration}
% =============================================================

\subsection{Certificate Chain Migration}

Migrating X.509 certificate chains to post-quantum signature algorithms
(ML-DSA, SLH-DSA, or FALCON) can proceed via two non-exclusive approaches:
sequential migration, in which a certificate authority issues a new
post-quantum-signed chain in parallel with the existing classical chain,
and composite certificates, in which a single certificate embeds both a
classical and a post-quantum public key and signature
\cite{x509_pqc_draft}. Composite certificates avoid the operational
overhead of maintaining two parallel chains but require relying-party
software capable of parsing the composite structure.

\subsection{JWT and OAuth 2.0 Signature Migration}

JWTs \cite{rfc7519} and the OAuth 2.0 framework \cite{rfc6749} are
signature-algorithm-agnostic by design, identifying the signing algorithm
in the JWT \texttt{alg} header. This permits incremental migration: an
identity provider may begin issuing ML-DSA-signed tokens while resource
servers retain the capacity to validate both classical (RS256/ES256) and
post-quantum signatures during a defined transition window. The principal
practical concern is token size: ML-DSA signatures (roughly 2.4 KB for
ML-DSA-44) are substantially larger than ECDSA signatures (64--72 bytes),
which has implications for HTTP header size limits when tokens are
transmitted as bearer tokens in the \texttt{Authorization} header.

\subsection{API Gateway Migration}

API gateways sit at the convergence point of transport, certificate, and
token layer migrations, and therefore require coordinated configuration
changes: enabling hybrid TLS key-exchange groups, trusting dual-algorithm
certificate chains, and validating tokens signed under either classical or
post-quantum algorithms. We recommend gateway configuration be
externalized (rather than hard-coded) to support algorithm rollback during
the transition, consistent with the cryptographic agility principle in
Section~\ref{sec:arch}.

% =============================================================
\section{Migration Framework}
% =============================================================

We propose a four-phase migration framework, detailed further in
\texttt{docs/migration-guide.md}:

\begin{enumerate}
  \item \textbf{Phase 0 -- Inventory:} Catalog all cryptographic
        dependencies across transport, certificate, token, and gateway
        layers, including third-party libraries and vendor-managed
        infrastructure not under direct organizational control.
  \item \textbf{Phase 1 -- Hybrid Transport:} Enable hybrid TLS key
        exchange at gateway and load-balancer tiers, defaulting to
        classical-only fallback for non-upgraded clients.
  \item \textbf{Phase 2 -- Certificate and Token Migration:} Issue
        composite or parallel post-quantum certificate chains; begin
        dual-algorithm JWT issuance and validation.
  \item \textbf{Phase 3 -- Agility Consolidation:} Remove deprecated
        classical-only code paths once client population and compliance
        requirements permit, retaining configuration-level agility for
        future algorithm transitions.
\end{enumerate}

This sequencing prioritizes transport-layer protection first, since
HNDL exposure (Section~\ref{sec:arch}) accrues continuously from the
present moment, while certificate and token migration -- which affect
authenticity rather than confidentiality -- can follow on a slightly
longer timeline without retroactive exposure.

% =============================================================
\section{Discussion and Open Challenges}
% =============================================================

Several engineering challenges remain open. First, the increased size of
post-quantum keys, ciphertexts, and signatures (Table~\ref{tab:algo-comparison})
compounds across layered protocols -- a TLS handshake, certificate chain,
and bearer token may each independently grow, with cumulative effects on
middlebox compatibility (e.g., devices that assume bounded TLS record
sizes) that are not yet fully characterized in production environments.
Second, composite certificate formats remain in active IETF
standardization \cite{x509_pqc_draft} and are not yet uniformly supported
across TLS libraries, creating near-term interoperability risk. Third,
maintaining dual classical/post-quantum trust chains during the transition
period introduces operational complexity -- key management, certificate
revocation, and audit processes must account for two parallel algorithm
families simultaneously. Finally, cryptographic agility, while a sound
design principle, is not free: configuration-driven algorithm selection
introduces additional testing surface and the risk of misconfiguration
(e.g., inadvertently permitting fallback to weak classical algorithms in
production).

This paper does not provide original empirical performance measurements;
the prototype, benchmark, and experimental evaluation components described
in the accompanying repository are scoped as future work, intended to
quantify the latency, throughput, and packet-size effects discussed
qualitatively above.

% =============================================================
\section{Conclusion}
% =============================================================

The transition to post-quantum cryptography in API systems is not a single
algorithm substitution but a multi-layer architectural migration spanning
transport, certificate, token, and gateway components. We have presented a
five-layer reference architecture, a threat model accounting for both
classical and harvest-now-decrypt-later adversaries, and a phased migration
framework that prioritizes hybrid transport-layer protection while
allowing certificate and token migration to proceed on a coordinated but
distinct timeline. Cryptographic agility -- the capacity to substitute
algorithms without re-architecting the system -- should be treated as a
design requirement throughout this process, given the likelihood of
continued evolution in post-quantum standards. Future work, outlined in
the accompanying repository's research roadmap, includes empirical
benchmarking of hybrid TLS handshake latency, ML-KEM and ML-DSA performance
under representative API workloads, and prototype implementation of the
gateway-layer migration patterns described in Section~\ref{sec:arch}.

% =============================================================
\section*{Acknowledgment}
% =============================================================

The author thanks the maintainers of the Open Quantum Safe project
\cite{liboqs} for their continued work providing accessible reference
implementations of NIST post-quantum algorithms, which inform the
prototype work planned in this repository's research roadmap.

% =============================================================
\bibliographystyle{IEEEtran}
\bibliography{references}

\end{document}
